\begin{textsample}{POS Dim 2 – 2014\_c5 – Score -1.00 – 2014\_c5/enade\_2014\_eng\_qui\_2e9817b8f4.txt}  \label{ex:f2_pos_314}
CAPÍTULO 5 DISTRIBUIÇÃO DOS CONCEITOS 5.1 PANORAMA NACIONAL DA DISTRIBUIÇÃO DOS CONCEITOS O Conceito Enade é um indicador de qualidade que avalia o desempenho dos estudantes a partir dos resultados do ENADE.
Os valores possíveis vão de 1 ( pior situação ) a 5 ( melhor situação ).
As Unidades de Observação com apenas um ou sem nenhum concluinte participante não obtêm o Conceito Enade, ficando Sem Conceito ( SC ).
A Tabela 5.1 apresenta a quantidade e distribuição de cursos de Engenharia Química participantes do ENADE/2014, por faixa de conceito e Grande Região.
Mais uma vez, ressaltando que a diferença entre os cursos tabulados neste capítulo e no capítulo 2 são os cursos sem conceito, em princípio, aqueles sem alunos concluintes que participassem da prova ou que tivessem tão somente um aluno concluinte21.
Observando -se os dados da Tabela 5.1, nota -se que, dos 102 cursos participantes, 43 ( 42,2 \% ) classificaram -se com conceito 3, o valor modal.
Este foi também o conceito modal em todas as demais regiões.
O conceito 2 foi o segundo mais frequente em nível nacional ( 29,4 \%, correspondendo a 30 cursos ) e o conceito 4, o terceiro ( 17,6 \%, correspondendo a 18 cursos ).
Houve, ainda, seis cursos ( 5,9 \% ) que receberam conceito 5 e quatro ( 3,9 \% ) cursos que receberam conceito 1.
Dos 102 cursos de Engenharia Química, apenas um curso ( 1,0 \% ) ficaram sem conceito ( SC ).
A região Norte participou com cinco cursos ou 4,9 \% do total nacional.
Nesta região, o conceito 3 foi atribuído a quatro cursos dos cinco participantes, o que equivale a 80,0 \% do total regional, além de ser o conceito modal na região.
Um curso ficou com conceito 1.
Nenhum curso ficou sem conceito ( SC ) ou com os demais conceitos restantes nesta região. 21 Estes cursos com somente um concluinte não foram considerados no capítulo 2. equivale a 80,0 \% do total regional, além de ser o conceito modal na região.
Um curso ficou com conceito 1.
Nenhum curso ficou sem conceito ( SC ) ou com os demais conceitos restantes nesta região.
A região Nordeste participou com 15 cursos ou 14,7 \% do total nacional.
Nesta região, seis cursos, 40,0 \% em termos regionais, obtiveram conceito 3, o conceito modal para a região.
Os conceitos 2 e 4 foram atribuídos a, respectivamente, quatro e três cursos ( 26,7 \% e 20,0 \% ).
Dois cursos obtiveram os conceitos 1 e 5 ( um curso, cada ), correspondendo a 6,7 \% cada.
Nenhum curso ficou sem conceito ( SC ).
Dos 55 cursos participantes da região Sudeste ( 53,9 \% do total nacional ), o conceito 3 também foi o conceito modal, e foi obtido por 20 ( 36,4 \% ) cursos.
O conceito 1 foi atribuído a um curso ( 1,8 \% ) e os conceitos 2 e 4, com 19 e 11 cursos, correspondendo a 34,5 \% e 20,0 \%, respectivamente.
O conceito 5 foi atribuído a quatro cursos ( 7,3 \% ).
E nenhum curso ficou sem conceito ( SC ).
Os 26 cursos da região Sul ficaram distribuídos em todas as faixas de conceitos.
Há predominância do conceito 3 com 46,2 \% ( conceito modal ), ou seja, 12 dos 26 cursos participantes na região Sul.
O conceito 2 foi atribuído a sete cursos ( 26,9 \% ), o conceito 1 a um curso ( 3,8 \% ) e o conceito 4, a quatro cursos ( 15,4 \% ).
O conceito 5 foi recebido por um curso ( 3,8 \% ).
Um curso ( 3,8 \% ) na região Sul ficou sem conceito ( SC ).
A região Centro-Oeste participou com somente um curso, e este curso foi avaliado com conceito 3.
O Gráfico 5.1 apresenta a distribuição do Conceito Enade segundo Grande Região : poligonais mais à esquerda representam Grandes Regiões com pior distribuição, e poligonais mais à direita, Grandes Regiões com melhores conceitos.
As regiões Centro-Oeste ( linha preta ) e Norte ( linha azul ) oferecem poucos cursos e as poligonais correspondentes são pouco suaves.
As regiões Sudeste ( linha vermelha ), Nordeste ( linha verde ) e Sul ( linha roxa ) apresentam distribuições bem semelhantes. 5.2 CONCEITOS POR CATEGORIA ADMINISTRATIVA E POR GRANDE REGIÃO A Tabela 5.2 apresenta a distribuição dos cursos participantes do ENADE/2014 de Engenharia Química, por Categoria Administrativa, de acordo com os conceitos por eles alcançados, segundo as Grandes Regiões brasileiras.
Dos 102 cursos participantes, 51 ( 50,0 \% ) eram ministrados em Instituições Públicas e 51 ( 50,0 \% ), em Privadas.
De acordo com as informações da Tabela 5.2, em termos nacionais, entre Instituições Públicas e Privadas : das seis IES que receberam o conceito 5, todas eram Públicas.
Dos 51 cursos participantes de IES Públicas, o conceito 3 foi o valor modal, atribuído a 20 cursos ( 39,2 \% ).
Entre os demais cursos de Instituições Públicas participantes, um curso ( 2,0 \% ) foi avaliado com conceito 1, oito cursos obtiveram conceito 2 ( 15,7 \% da categoria ), 15 cursos receberam conceito 4 ( 29,4 \% ) e, como já comentado, seis foram alocados ao conceito 5 ( 11,8 \% ).
Nesta categoria, somente um curso ( 2,0 \% ) ficou sem conceito ( SC ).
Na rede privada, o conceito modal foi 3, com 23 cursos ( 45,1 \% ) dos 51 da categoria.
Entre os demais cursos participantes, três receberam conceito 1 ( 5,9 \% ), e 22, conceito 2 ( 43,1 \% ).
O conceito 4 foi atribuído a três cursos ( 5,9 \% ).
Nenhum curso ficou sem conceito ( SC ) ou obteve conceito 5.
Na análise por região, observa -se que, na região Norte, as Instituições Privadas participaram com somente um curso ( 20,0 \% do total regional ), e este curso foi avaliado com conceito 1.
As Instituições Públicas participaram com quatro cursos na região Norte ( 80,0 \% do total regional ).
Todos os foram avaliados com conceito 3.
Na região Nordeste, a rede privada concentrou quatro dos 15 cursos participantes, o equivalente a 26,7 \% do total da região.
Dentre os cursos oferecidos por IES Privadas no Nordeste, um curso ficou alocado ao conceito 1, outro curso obteve conceito 3 e dois cursos ficaram com conceito 2, conceito modal.
Nenhum curso ficou sem conceito ( SC ) ou obteve conceitos 4 ou 5.
As Instituições Públicas dessa região participaram com 11 cursos ( 73,3 \% ), dos quais cinco ( 45,5 \% ) obtiveram conceito 3, o valor modal, seguido por três cursos que receberam conceito 4 ( 27,3 \% ).
Dois cursos foram alocados ao conceito 2 ( 18,2 \% ), enquanto nenhum curso ficou sem conceito ( SC ) ou com conceito 1.
Somente um curso ( 9,1 \% ) foi avaliado com conceito 5.
Na região Sudeste, a proporção de cursos da rede privada, 63,6 \%, foi mais elevada do que nas demais regiões brasileiras, correspondendo a 35 dos 55 cursos participantes.
Nesta categoria, na região Sudeste, o conceito modal foi 2 ( 17 cursos ).
Nesta combinação de Categoria Administrativa e Grande Região, nenhum curso não recebeu conceito ( SC ) ou obteve conceito 5.
Os demais foram avaliados com conceito 1 ( um curso ), conceito 3 ( 14 cursos ) e conceito 4 ( três cursos ).
Entre os 20 cursos oferecidos por Instituições Públicas na região Sudeste ( 36,4 \% ), menor proporção desta categoria dentre as regiões, o conceito modal foi 4 ( oito cursos ), e os demais cursos receberam os conceitos 2 ( dois cursos ), 3 ( seis cursos ) e 5 ( quatro cursos ).
Na região Sudeste, nesta categoria, nenhum curso ficou sem conceito ( SC ) ou recebeu conceito 1.
As Instituições Privadas concentraram 11 dos 26 cursos participantes da região Sul, 42,3 \% do total regional.
Desses, oito ficaram com conceito 3 ( o conceito modal ) e três com conceito 2.
Nenhum curso ficou sem conceito ( SC ) ou recebeu os demais conceitos restantes nesta categoria e região.
As Instituições Públicas na região Sul participaram com 15 cursos ( 57,7 \% ) e os conceitos modais foram 2, 3 e 4 ( quatro cursos, cada ).
Dois cursos ficaram com conceitos 1 e 5 ( um curso, cada ).
Somente um curso ficou sem conceito ( SC ) nesta categoria e região.
Na região Centro-Oeste, as Instituições Privadas não tiveram participação.
As Instituições Públicas participaram com somente um curso e este curso foi avaliado com conceito 3.
O Gráfico 5.2 apresenta a distribuição do Conceito Enade segundo a Categoria Administrativa da IES.
Os cursos em IES Públicas ( linha azul ) apresentam uma poligonal mais à direita do que os em IES Privadas ( linha verde ) e, consequentemente, uma distribuição de Conceitos Enade com valores maiores. 5.3 CONCEITOS POR ORGANIZAÇÃO ACADÊMICA E POR GRANDE REGIÃO Na Tabela 5.3, encontra -se a distribuição dos conceitos atribuídos aos cursos participantes do ENADE/2014 na Área de Engenharia Química, por Organização Acadêmica, segundo as Grandes Regiões brasileiras.
Dos 102 cursos de Engenharia Química participantes, 71 eram oferecidos em Universidades, 18 em Faculdades e 13 em Centros Universitários.
Esta distribuição corresponde a, respectivamente, 69,6 \%, 17,6 \% e 12,7 \% dos cursos.
De acordo com os dados apresentados, de todos os seis cursos avaliados com conceito 5, cinco eram vinculados a Universidades.
Somente um curso oferecido em Universidades ficou sem conceito ( SC ).
Esse tipo de Organização Acadêmica teve o conceito 3 como modal, com 29 cursos ( 40,8 \% ).
Os demais cursos de Universidades avaliados receberam os conceitos 1 ( um curso ), 2 ( 18 cursos ), 4 ( 17 cursos ) e conceito 5 ( cinco cursos, como já mencionado ).
Entre os cursos em Centros Universitários, o conceito modal foi 3, com seis cursos ( 46,2 \% ).
Neste tipo de Organização Acadêmica, nenhum curso ficou sem conceito ( SC ) ou recebeu os conceitos 4 ou 5.
Cinco cursos receberam o conceito 2 e dois, conceito 1.
Nas Faculdades, nenhum dos 18 cursos ficou sem conceito ( SC ), e oito ( 44,4 \% ) ficaram alocados no conceito 3, conceito modal, seguido de sete cursos que obtiveram conceito 2 e de três cursos que receberam os conceitos 1, 4 e 5 ( um curso, cada ).
Considerando -se separadamente as regiões brasileiras, verifica -se que, na região Norte, as Universidades concentraram quatro dos cinco cursos participantes, todos avaliados com conceito 3.
Os Centros Universitários da região Norte foram representados por somente um curso, e este recebeu o conceito 1.
As Faculdades não tiveram participação na região.
Na região Nordeste, Universidades concentraram 14 dos 15 cursos da região.
Entre os cursos em Universidades, nessa região, o conceito modal foi 3 com seis cursos, enquanto nenhum curso ficou sem conceito ( SC ) ou com conceito 1.
Os demais cursos receberam os conceitos 2 ( quatro cursos ), 4 ( três cursos ) e 5 ( um curso ).
Os Centros Universitários não tiveram participação na região Nordeste.
As Faculdades participaram com somente um curso na região Nordeste, e este curso obteve conceito 1.
Na região Sudeste, as Universidades concentraram 28 dos 55 cursos da região.
Entre os cursos em Universidades, nessa região, o conceito modal foi 4 com dez cursos, enquanto nenhum curso ficou sem conceito ( SC ) ou com conceito 1.
Os demais cursos receberam os conceitos 2 ( sete cursos ), 3 ( oito cursos ) e 5 ( três cursos ).
Os Centros Universitários participaram com 11 cursos na região Sudeste, dos quais dez obtiveram os conceitos modais, 2 e 3 ( cinco cursos, cada ).
Um curso ficou com conceito 1.
Nenhum curso ficou sem conceito ( SC ) ou com conceitos 4 ou 5.
As Faculdades foram representadas por 16 cursos na região Sudeste dos quais 14 receberam os conceitos 2 e 3 ( sete cursos, cada ).
Os demais se distribuíram nos conceitos 4 ( um curso ) 5 ( um curso ).
Nenhum curso ficou sem conceito ( SC ) ou com conceito 1.
Dos 26 cursos da região Sul, 24 eram de Universidades, para os quais o conceito modal foi 3, com dez cursos.
Nesse tipo de organização, um dos cursos ficou sem conceito ( SC ).
Os demais receberam os conceitos 1 ( um curso ), 2 ( sete cursos ), 4 ( quatro cursos ) e 5 ( um curso ).
Os Centros Universitários da região Sul participaram com somente um curso, e este curso foi avaliado com conceito 3.
As Faculdades na região Sul também participaram com apenas um curso, e este curso foi obteve conceito 3.
Na região Centro-Oeste, um curso era de Universidades, e este curso obteve conceito 3.
Os Centros Universitários e as Faculdades da região Centro-Oeste não tiveram participação.
O Gráfico 5.3 apresenta a distribuição do Conceito Enade segundo a Organização Acadêmica da IES.
Os cursos em Universidades ( linha azul ) apresentam uma poligonal mais à direita do que as demais, denotando uma melhor distribuição de Conceitos Enade.
No outro extremo, temos os cursos oferecidos em Centros Universitários ( linha verde ) com a poligonal mais à esquerda e, consequentemente, uma distribuição de Conceitos Enade com menores valores.
IES de Faculdades ( linha vermelha ) apresentam uma situação intermediária.

% matched lemmas: 
\end{textsample}
